\documentclass[11pt]{article}

\usepackage[margin=1in]{geometry}
\usepackage{amsmath,amssymb,amsthm,mathtools}
\usepackage{booktabs}
\usepackage{microtype}
\usepackage{enumitem}
\usepackage{hyperref}

\hypersetup{colorlinks=true,linkcolor=blue,citecolor=blue,urlcolor=blue}

\newtheorem{theorem}{Theorem}[section]
\newtheorem{proposition}[theorem]{Proposition}
\newtheorem{lemma}[theorem]{Lemma}
\newtheorem{corollary}[theorem]{Corollary}
\newtheorem{example}[theorem]{Example}
\theoremstyle{definition}
\newtheorem{definition}[theorem]{Definition}
\theoremstyle{remark}
\newtheorem{remark}[theorem]{Remark}

\newcommand{\ZZ}{\mathbb Z}
\newcommand{\CC}{\mathbb C}
\newcommand{\one}{\mathbf 1}
\newcommand{\Dd}{\mathcal D}
\newcommand{\QQ}{\mathcal Q}
\newcommand{\ee}{\mathrm e}
\newcommand{\Harm}{\mathsf H}
\newcommand{\Ecl}{\mathcal E_{\mathrm{cl}}}
\newcommand{\Edir}{\mathcal E_{\mathrm{dir}}}

\title{\textbf{Short-Quotient M\"obius Tails and a No-Power-Loss\\
Majorant for the Physical Cluster Gram}}
\author{Liang Wang\\
Huazhong University of Science and Technology\\
Wuhan 430074, P.R. China}
\date{August 2026}

\begin{document}
\maketitle

\begin{abstract}
The transition-native twin-prime program contains a squarefree divisor band
and, after complete-period frequency clustering, the M\"obius-log tails
\(
C_h=\sum_{h\mid d}\mu(d)\log(d)/d.
\)
This paper estimates those tails on the literal V46 scales.  The nonzero
integer cutoff in the reciprocal emitter forces every active reduced
denominator to satisfy
\(h\ge H/q_{\max}\ge H/(2Q)=2Y_0\).
Thus $h$ itself lies in the full transition band.  Writing $d=hk$ gives an
exact quotient M\"obius sum whose length is at most
\(2UQ/H=2x^{23/2400+o(1)}\).  The $k=1$ term anchors the direct coefficient
mass, while an elementary harmonic estimate bounds every cluster coefficient.
Together with an exact reduced-denominator decomposition of the emitter row
norm, this proves that the complete-period cluster Gram is at most
\(O((\log x)^2)=x^{o(1)}\) times the divisor direct-sum energy.  Hence
frequency clustering creates no fixed-power amplification.  The comparison
is not a saving theorem: every active denominator in $U/2<h\le U$ has
cluster-to-direct coefficient ratio exactly one.  The direct-sum arithmetic
energy, finite-window off-frequency Gram, prime-shell reassembly, arithmetic
$L^2$, and twin-prime endpoint remain open.
\end{abstract}

\paragraph{Claim level.}
\texttt{PROVED\_STRUCTURAL\_L1 / SHORT\_QUOTIENT\_CLUSTER\_MAJORANT}.\\
The $O((\log x)^2)$ comparison is proved on the source-locked asymptotic
scales.  Finite decimal ratios are numerical observations only.  No
arithmetic saving or twin-prime conclusion is claimed.

\section{Introduction}

The V46 transition scalar in the analytic twin-prime route contains one
reciprocal emitter for each squarefree divisor $d$ in a short proper-factor
range.  Its literal coefficient is
\begin{equation}
 c_d=\frac{\mu(d)\log d}{d}.
 \label{eq:coefficient}
\end{equation}
TPC-213 showed that these rows are pulled back from one common physical source,
so different divisor blocks are not an orthogonal direct sum
\cite{WangTPC213}.  TPC-214 then identified the exact complete-period
frequency clusters \cite{WangTPC214}.  If $h\mid d$, the row at a frequency
with reduced denominator $h$ equals the row at $h$ itself.  Consequently all
multiples of $h$ combine before the Hermitian square, with coefficient
\begin{equation}
 C_h=\sum_{\substack{d\in\Dd_x\\h\mid d}}c_d.
 \label{eq:cluster-coefficient}
\end{equation}

The remaining issue is quantitative.  A completely arbitrary family of
multiples could make the ratio between $|C_h|^2$ and
\(\sum_{h\mid d}|c_d|^2\) as large as the family size.  The literal V46
emitter, however, has an integer activation threshold.  A reduced row cannot
appear until its denominator is large enough to admit a nonzero reciprocal
frequency.  That threshold places $h$ back inside the original divisor band
and makes the quotient $k=d/h$ very short.

This observation leads to four exact steps:
\begin{enumerate}[leftmargin=2em]
 \item every active denominator satisfies $h\ge H/q_{\max}\ge2Y_0$;
 \item its cluster coefficient is a quotient M\"obius sum with
       $k\le2x^{23/2400+o(1)}$;
 \item the $k=1$ coefficient anchors a harmonic comparison with the direct
       coefficient mass;
 \item reduced fractions partition every divisor row, converting the
       coefficient comparison into a complete-period energy comparison.
\end{enumerate}

The resulting bound is subpower, and in fact $O((\log x)^2)$.  This is a real
structural advance over leaving the cluster tails unbounded, but it deliberately
stops before arithmetic promotion.  A top-shell identity also shows why the
argument cannot by itself produce a fixed-power saving.

\section{Literal source lock and the emitter}

We use the V46 parameters exactly as frozen in the transition-native Euler
artifact \cite{WangV46}:
\begin{equation}
 H=x^{21/32},\qquad Q=x^{1/3},\qquad
 Y_0=\frac{H}{4Q}=x^{31/96+o(1)},\qquad
 U=x^{133/400}.
 \label{eq:scales}
\end{equation}
Let
\begin{equation}
 \QQ_x=\{q\text{ prime}:Q<q\le2Q\},
 \qquad q_{\max}=\max\QQ_x,
 \label{eq:prime-shell}
\end{equation}
and let the full squarefree transition band be
\begin{equation}
 \Dd_x=\{d\in\ZZ:Y_0<d\le U,\ \mu(d)^2=1\}.
 \label{eq:divisor-band}
\end{equation}
For sufficiently large $x$, the source inequalities give $d<U<Q<q$.  Thus
every $q\in\QQ_x$ is invertible modulo every $d\in\Dd_x$.

For an arbitrary profile $\psi$ on the arguments below, define the
coefficient-free reciprocal row
\begin{equation}
 B_d(r)=\sum_{q\in\QQ_x}
 \sum_{0<|m|\le dq/H}
 \psi\!\left(\frac{Hm}{dq}\right)
 \one_{m\overline q\equiv r\pmod d},
 \qquad r\in\ZZ/d\ZZ.
 \label{eq:emitter}
\end{equation}
The cutoff is interpreted in the integer sense:
\(|m|\le\lfloor dq/H\rfloor\).

TPC-214 proves the exact dilation covariance
\begin{equation}
 B_d\!\left(\frac dh a\right)=B_h(a),
 \qquad h\mid d,
 \label{eq:dilation}
\end{equation}
and the complete-period cluster factorization.  We use those two identities
as the parent theorem; no continuous cutoff or independent-profile
substitution is introduced here.

\begin{remark}[the full-band hypothesis]
The family in \eqref{eq:divisor-band} contains every squarefree integer in the
stated interval.  This is essential below.  For an arbitrary subfamily, an
active reduced denominator can divide a selected $d$ without the diagonal
member $d=h$ being selected.
\end{remark}

\section{Activation and the short quotient}

\begin{lemma}[activation floor]
\label{lem:activation}
Suppose $h$ divides a member of $\Dd_x$ and $B_h$ is not the zero row.  Then
\begin{equation}
 h\ge\frac{H}{q_{\max}}
 \ge\frac{H}{2Q}=2Y_0.
 \label{eq:activation}
\end{equation}
In particular, $h\in\Dd_x$.
\end{lemma}

\begin{proof}
If $B_h$ is nonzero, then at least one contributing term in
\eqref{eq:emitter} has a nonzero integer $m$.  For the associated $q$,
\[
 1\le |m|\le\frac{hq}{H}\le\frac{hq_{\max}}{H}.
\]
This proves the first inequality in \eqref{eq:activation}; the second follows
from $q_{\max}\le2Q$ and the definition of $Y_0$.  Since $h$ divides a
squarefree $d\in\Dd_x$, it is squarefree and $h\le d\le U$.  Equation
\eqref{eq:activation} gives $h>Y_0$, so $h$ belongs to the full band.
\end{proof}

\begin{lemma}[short-quotient normal form]
\label{lem:normal-form}
For every active reduced denominator $h$,
\begin{equation}
 C_h=\frac{\mu(h)}{h}
 \sum_{\substack{Y_0/h<k\le U/h\\(k,h)=1}}
 \frac{\mu(k)}{k}\bigl(\log h+\log k\bigr).
 \label{eq:normal-form}
\end{equation}
The term $k=1$ occurs, and every nonzero summand satisfies
\begin{equation}
 k\le\frac{Uq_{\max}}{H}
 \le\frac{2UQ}{H}
 =2x^{23/2400+o(1)}.
 \label{eq:quotient-bound}
\end{equation}
\end{lemma}

\begin{proof}
Write each multiple of $h$ as $d=hk$.  Since $d$ is squarefree,
$(h,k)=1$ and $\mu(hk)=\mu(h)\mu(k)$.  Substitution into
\eqref{eq:cluster-coefficient} gives \eqref{eq:normal-form}; allowing
nonsquarefree $k$ is harmless because their M\"obius coefficient is zero.
Lemma~\ref{lem:activation} shows that $d=h$ belongs to $\Dd_x$, so $k=1$
is present.  It also gives
\[
 k\le\frac Uh\le\frac{Uq_{\max}}H\le\frac{2UQ}H.
\]
Finally,
\[
 \frac{133}{400}+\frac13-\frac{21}{32}
 =\frac{798+800-1575}{2400}
 =\frac{23}{2400},
\]
which proves the last equality in \eqref{eq:quotient-bound}.
\end{proof}

\begin{remark}[scale of the reduction]
The active divisor itself can be as large as $x^{133/400}$, but after fixing
an emitter-visible reduced denominator, all cluster multiples are indexed by
a quotient of length only $x^{23/2400+o(1)}$.  This is a deterministic support
reduction, not M\"obius cancellation.
\end{remark}

\section{A deterministic coefficient majorant}

For every active $h$, define the direct coefficient mass
\begin{equation}
 D_h=\sum_{\substack{d\in\Dd_x\\h\mid d}}|c_d|^2.
 \label{eq:direct-mass}
\end{equation}
Let $\Harm_n=\sum_{k=1}^n1/k$ denote the $n$th harmonic number, with
$\Harm_0=0$.

\begin{lemma}[anchored harmonic comparison]
\label{lem:harmonic}
For every active $h$,
\begin{equation}
 |C_h|^2\le
 \left(\frac{\log U}{\log h}\right)^2
 \Harm_{\lfloor U/h\rfloor}^{,2}D_h.
 \label{eq:row-majorant}
\end{equation}
\end{lemma}

\begin{proof}
The $d=h$ term supplied by Lemma~\ref{lem:activation} gives
\begin{equation}
 D_h\ge |c_h|^2=\left(\frac{\log h}{h}\right)^2.
 \label{eq:anchor}
\end{equation}
On the other hand, the triangle inequality and $d=hk\le U$ give
\begin{align}
 |C_h|
 &\le \sum_{\substack{k\le U/h\\(k,h)=1}}
       \frac{|\mu(k)|\log(hk)}{hk} \\
 &\le \frac{\log U}{h}
       \sum_{k\le U/h}\frac1k
 =\frac{\log U}{h}\Harm_{\lfloor U/h\rfloor}.
 \label{eq:triangle}
\end{align}
Squaring \eqref{eq:triangle} and using \eqref{eq:anchor} proves
\eqref{eq:row-majorant}.
\end{proof}

Define the source-locked uniform factor
\begin{equation}
 A_x=\left(\frac{\log U}{\log(H/q_{\max})}\right)^2
 \Harm_{\lfloor Uq_{\max}/H\rfloor}^{\,2}.
 \label{eq:Ax}
\end{equation}

\begin{corollary}[subpower coefficient comparison]
\label{cor:subpower}
Uniformly over all active reduced denominators,
\begin{equation}
 |C_h|^2\le A_xD_h,
 \qquad A_x=O((\log x)^2)=x^{o(1)}.
 \label{eq:subpower-coefficient}
\end{equation}
More explicitly,
\begin{equation}
 A_x\le
 \left(\frac{\log U}{\log(H/(2Q))}\right)^2
 \Harm_{\lfloor2UQ/H\rfloor}^{,2}.
 \label{eq:explicit-Ax}
\end{equation}
\end{corollary}

\begin{proof}
Both factors on the right of \eqref{eq:row-majorant} are nonincreasing as $h$
increases.  Apply Lemma~\ref{lem:activation} to obtain
\eqref{eq:subpower-coefficient}.  Since $q_{\max}\le2Q$, enlarging the
harmonic index and decreasing the logarithmic denominator gives
\eqref{eq:explicit-Ax}.  The logarithm ratio is bounded on the scales
\eqref{eq:scales}, while
\[
 \Harm_n\le1+\log n,
 \qquad n\le2x^{23/2400+o(1)}.
\]
Thus $A_x=O((\log x)^2)$, which is $x^{o(1)}$.
\end{proof}

\section{From rows to the complete-period Gram}

For $h\ge2$, set
\begin{equation}
 N_h=\sum_{\substack{a\bmod h\\(a,h)=1}}|B_h(a)|^2.
 \label{eq:primitive-norm}
\end{equation}
For completeness, let $N_1=|B_1(0)|^2$.  In the source range $q<H$, so the
additive zero axis vanishes and $N_1=0$ \cite{WangTPC214}.

\begin{lemma}[row-norm divisor decomposition]
\label{lem:row-decomposition}
For every $d\in\Dd_x$,
\begin{equation}
 \sum_{r\bmod d}|B_d(r)|^2=\sum_{h\mid d}N_h.
 \label{eq:row-decomposition}
\end{equation}
\end{lemma}

\begin{proof}
Each residue $r\bmod d$ determines a unique reduced fraction
\[
 \frac rd=\frac ah,
 \qquad h\mid d,\quad (a,h)=1,
\]
where the zero fraction is represented by $a=0,h=1$.  Conversely, each pair
on the right represents the residue $r=(d/h)a$.  The classes indexed by $h$
therefore partition the complete residue row.  Dilation covariance
\eqref{eq:dilation} gives $B_d(r)=B_h(a)$ in each class, proving
\eqref{eq:row-decomposition}.
\end{proof}

Let
\begin{equation}
 L=\operatorname{lcm}\{d:d\in\Dd_x\}
 \label{eq:period}
\end{equation}
for a finite realization of the transition band, and define
\begin{align}
 \Ecl
 &=\sum_{u\bmod L}\left|
   \sum_{d\in\Dd_x}c_d\sum_{r\bmod d}
   B_d(r)\ee^{2\pi iru/d}\right|^2,
 \label{eq:cluster-energy}\\
 \Edir
 &=L\sum_{d\in\Dd_x}|c_d|^2
   \sum_{r\bmod d}|B_d(r)|^2.
 \label{eq:direct-energy}
\end{align}

\begin{theorem}[no-power-loss cluster majorant]
\label{thm:main}
On the literal V46 squarefree transition band,
\begin{equation}
 \boxed{\Ecl\le A_x\Edir},
 \qquad A_x=O((\log x)^2)=x^{o(1)},
 \label{eq:main}
\end{equation}
where $A_x$ is given explicitly by \eqref{eq:Ax} and
\eqref{eq:explicit-Ax}.
\end{theorem}

\begin{proof}
The TPC-214 complete-period factorization gives
\begin{equation}
 \frac{\Ecl}{L}=\sum_hN_h|C_h|^2.
 \label{eq:tpc214-factor}
\end{equation}
Corollary~\ref{cor:subpower} and a finite rearrangement of nonnegative terms
give
\begin{align}
 \frac{\Ecl}{L}
 &\le A_x\sum_hN_hD_h \\
 &=A_x\sum_{d\in\Dd_x}|c_d|^2\sum_{h\mid d}N_h \\
 &=A_x\sum_{d\in\Dd_x}|c_d|^2
       \sum_{r\bmod d}|B_d(r)|^2
 =A_x\frac{\Edir}{L},
\end{align}
where the penultimate equality is Lemma~\ref{lem:row-decomposition}.
Multiplication by $L$ proves \eqref{eq:main}.
\end{proof}

\begin{remark}[what the theorem does and does not remove]
The theorem proves that the complete-period shared-frequency Gram has no
additional fixed-power cost beyond the divisor direct sum.  It does not bound
\(\Edir\) at the natural arithmetic scale.  It also does not replace the
finite physical interval by a complete period: off-frequency terms that
vanish in \eqref{eq:tpc214-factor} must be revisited in the actual window.
\end{remark}

\section{A sharp top-shell obstruction}

\begin{proposition}[exact ratio one in the top shell]
\label{prop:top-shell}
If an active reduced denominator satisfies $U/2<h\le U$, then
\begin{equation}
 C_h=c_h,\qquad D_h=|c_h|^2,
 \qquad \frac{|C_h|^2}{D_h}=1.
 \label{eq:top-shell}
\end{equation}
\end{proposition}

\begin{proof}
Lemma~\ref{lem:activation} places $h$ itself in $\Dd_x$.  Any distinct
positive multiple of $h$ is at least $2h>U$, so no other member of the band is
divisible by $h$.  The definitions \eqref{eq:cluster-coefficient} and
\eqref{eq:direct-mass} reduce to the single term $d=h$.
\end{proof}

Proposition~\ref{prop:top-shell} is a sharp scoped obstruction.  It refutes a
uniform rowwise estimate of the form
\begin{equation}
 |C_h|^2\le x^{-\delta}D_h
 \qquad(\delta>0)
 \label{eq:forbidden-saving}
\end{equation}
based only on cluster coefficients.  It does not refute a global saving that
uses the size or distribution of $N_h$, finite-window oscillation, prime-shell
averaging, or the signed four-packet scalar.  Those inputs are absent from the
coefficient identity \eqref{eq:top-shell}.

\section{Finite exact certificate}

The release fixture retains the TPC-214 rational profile but enlarges the
divisor family to a full finite band:
\begin{equation}
 \QQ=\{11,13,17\},\quad H=40,\quad Y_0=2,\quad U=35,
 \quad \psi(t)=\frac1{(1+t^2)^2}.
 \label{eq:fixture}
\end{equation}
The finite band consists of all squarefree $2<d\le35$ coprime to the three
fixture primes.  This coprimality, rather than the asymptotic inequality
$U<Q$, makes every finite inverse legal.  All emitter entries and row-norm
identities are exact rational numbers.

The activation floor is
\(
\lceil40/17\rceil=3
\), and the active denominators are
\begin{equation}
 3,5,6,7,10,14,15,19,21,23,29,30,31,35.
 \label{eq:active-fixture}
\end{equation}
The largest realized quotient is $10$, below the uniform integer bound
\(\lfloor35\cdot17/40\rfloor=14\).  Selected coefficient ratios are shown in
Table~\ref{tab:fixture}.  They use real logarithms and are therefore labeled
\texttt{NUMERICAL\_OBSERVATION}; the top-shell membership and ratio-one
statements are exact.

\begin{table}[t]
\centering
\small
\begin{tabular}{@{}rllrr@{}}
\toprule
$h$ & band multiples & quotients & $|C_h|^2/D_h$ & harmonic majorant \\
\midrule
3  & $3,6,15,21,30$ & $1,2,5,7,10$ & $0.0721263$ & $95.5111$ \\
5  & $5,10,15,30,35$ & $1,2,3,6,7$ & $0.0279988$ & $32.8075$ \\
6  & $6,30$ & $1,5$ & $0.3363555$ & $20.5278$ \\
7  & $7,14,21,35$ & $1,2,3,5$ & $0.1711627$ & $17.4043$ \\
15 & $15,30$ & $1,2$ & $0.0992569$ & $3.8782$ \\
19 & $19$ & $1$ & $1$ & $1.4580$ \\
35 & $35$ & $1$ & $1$ & $1$ \\
\bottomrule
\end{tabular}
\caption{Finite coefficient diagnostics.  The majorant is deliberately
cancellation-free.  Rows $h=19,21,23,29,30,31,35$ lie in the top shell and
have exact ratio one.}
\label{tab:fixture}
\end{table}

Exact dilation covariance and the row decomposition are checked for every
divisor and every reduced residue.  With the exact rational row norms followed
by real-logarithm evaluation, the whole finite cluster and direct energies are
\begin{equation}
 \Ecl/L\approx5.5554404644,qquad
 \Edir/L\approx9.3063240426,qquad
 \Ecl/\Edir\approx0.5969532588.
 \label{eq:finite-ratio}
\end{equation}
Equation \eqref{eq:finite-ratio} is a reproduction diagnostic, not evidence
for asymptotic cancellation.  An independent implementation reconstructs the
rows, formal log polynomials, harmonic factors, and top-shell list without
importing the producer.  Normal and optimized interpreter modes are required
to give byte-identical output.  A separate adversarial checker covers 53
active rows across three finite configurations and verifies that deleting the
$d=h$ member invalidates the diagonal-anchor hypothesis.

\section{Route evaluation and open arithmetic}

The main theorem changes the structural state in one precise way:
\begin{equation}
 \text{complete-period cluster energy}
 \ \le\ x^{o(1)}\times
 \text{divisor direct-sum row energy}.
 \label{eq:route-change}
\end{equation}
Thus a future proof no longer needs to budget a fixed power for coherent
reduced-frequency collisions.  The next arithmetic object is the right-hand
side of \eqref{eq:route-change}, not the unestimated coefficient tail.

Three independent gaps remain:
\begin{enumerate}[leftmargin=2em]
 \item \textbf{Direct-sum arithmetic energy.}  The current argument gives no
       natural-scale estimate for the physical residual rows weighted by
       $|c_d|^2$.
 \item \textbf{Finite physical interval.}  Complete-period orthogonality
       removes unequal rational frequencies.  The actual interval produces an
       off-frequency Gram that needs its own bound.
 \item \textbf{Prime-shell and four-packet reassembly.}  The prime-only shell,
       block localization, exact diagonal subtraction, and signed
       polarization must be retained in one theorem with the strict
       $1/400$ loss ledger.
\end{enumerate}

Accordingly, the formal Route-B status is
\begin{equation*}
\begin{array}{ll}
\texttt{STRUCTURAL\_THRESHOLD\_A} &= \texttt{PASS},\\
\texttt{CLUSTER\_TO\_DIRECT\_MAJORANT} &=
  \texttt{PROVED\_O\_LOG\_X\_SQUARED},\\
\texttt{UNIFORM\_ROWWISE\_POWER\_SAVING} &=
  \texttt{REFUTED\_SCOPED},\\
\texttt{ARITHMETIC\_ADVANCE} &= \texttt{NO},\\
\texttt{L2} &= \texttt{NONE},\\
\texttt{FULL\_GATE\_B\_STRICT\_1\_OVER\_400} &= \texttt{UNPAID}.
\end{array}
\end{equation*}
Route A is not applicable to this analytic twin-prime session.  No fixed-atom
credit or twin-prime conclusion follows from a complete-period comparison.

\section{Conclusion}

The literal reciprocal cutoff turns every emitter-visible M\"obius cluster
into a short quotient sum.  Because activation places the reduced denominator
itself back in the full divisor band, the diagonal term anchors a deterministic
harmonic comparison.  Reduced-fraction row decomposition then upgrades this
coefficient estimate to an $O((\log x)^2)$ complete-period Gram majorant.

This closes the fixed-power amplification question at the cluster stage.  The
top shell simultaneously proves that no uniform rowwise saving can come from
cluster algebra alone.  The next useful step is therefore to estimate the
physical direct-sum row energy before reintroducing finite-window cross
frequencies and the prime-shell four-packet assembly.

\bibliographystyle{plain}
\bibliography{references}

\end{document}
